On February 2, four of us set out from Tuba City. We were a selected group. We went to Winslow. At dusk we took off in a plane. We landed in Phoenix. It took us 30 minutes to get there.

Then early on the following morning we were awakened, and taken off in a car. We went to the Navy Field. At six o'clock we were put aboard a plane. There were four planes. We each got aboard a different one. I was sent aboard the plane called ``Tomahawk.'' It was a four-motored plane. Planes of this type can carry a huge load. I found a great pile of alfalfa aboard. These planes can even carry automobiles.

Then we took off for the Navajo country. We flew over Winslow. When we passed the reservation boundary they said to me, ``From now on it's up to you. You know where the people live.'' Then we flew over the lower end of the black lava peaks, and over Oraibi. And then over the lower end of the mountains and over Cow Springs. And then on over Shonto and over the ridge. And then on down to Oljato. Then we flew over Mexican Hat. While we were flying through there, there were other planes flying alongside us. They remained at the same distance apart. We had earphones on, and talked to one another by radio. This way we discussed our position with one another. This is how we arrived. And then we flew up along the San Juan River. Before we got to Shiprock we turned to go alongside the mountain range, toward Lukachukai. We circled several times over Lukachukai, and then we went toward Chinle. And then over Ganado and Klagetoh, to the railroad. Then we turned around and went back over Greasewood and Sunrise Springs. Then we followed the ridge up over Black Mountain. We continued to Kayenta where we turned around. And then we started back to Phoenix, where we arrived at noon. We ate and then flew right back over on the same route, and everything that had happened seemed like a dream. It was late in the afternoon when we got back to Phoenix. That's how we went about over our land dropping hay. In the late evening we all returned to Phoenix. Like this I went across our land unloading bales of hay. And they told us there'd be another flight on the next morning.

On the next day, February 4, we took off again. We were awakened from sleep at four o'clock, before dawn. However, we didn't take off until six o'clock. That day the fog was really thick. You can't tell where you are in that kind of weather. On that day I was up 18,000 feet in the sky -- I don't know how far above the clouds. I kept asking how high we were. And the pilot would tell me we were 18,000 feet above the ground (i.e. above sea level). I was just a little bit scared. The plane would repeatedly drop (in an air pocket). It would suddenly drop as much as 500 feet. The pilot told me he didn't blame me for being scared. It's a life or death matter when you ride a plane. Nevertheless, we hauled two loads that day. Nonetheless we hauled two loads again by the time the sun went down.

On the next day, which was February 5, we met with the Governor of Arizona. We also met with an army officer from California. These two men got aboard and we took off again. They radioed ahead that these men were aboard such and such a plane. We were loaded with hay, food and firewood. We dumped these on the Hopis. We dropped them on Walpi, Shipaulovi, Shimopovi, Oraibi, Hotevilla, and Ndadziizdeli. We made about two runs over these villages. As we flew over them we watched and saw people swarming about. We said they must be fighting over the firewood.

These men with whom we flew asked me many questions. ``To what extent are the people suffering,'' they asked me. I pointed out to them the various things the people were suffering from. After we had inspected the reservation we flew back, but to Tucson. There the two men got off. As for us, we flew back and got off at Phoenix.

On the next day, which was the sixth, I didn't go out, because the plane had made me somewhat deaf. That happened to two of us.

On Monday the seventh I went again. We went out according to our instructions. That day we made four flights, hauling only hay.

On the eighth we again went out. The main thing we hauled was hay. Like on that other occasion, there was again a lot of fog. It looked strange up above the clouds. In some places the clouds looked like mountain ranges; some looked like sand dunes, and some exactly like mesas. You couldn't see the earth below. Nevertheless, we carried four loads. In the evening, on our return trip, one of the motors caught fire. It happened between Flagstaff and Phoenix. ``Look what happened?'' ``It's burning!'' said the pilot. I looked and saw a fire sitting out alongside. There's nothing you can do when such things happen. It's a scary situation. You never can tell what will happen. Even if you jump out you'll lose your life. Travel by plane means that death is involved. It cannot be known. You just sit and hope for the best. I racked my brains trying to think how we were going to put out the fire. But there is a fire extinguisher right alongside these motors so when something like this happens the fire is easily put out.

When we got back we were told that we could have the next day off. This was welcome news. It was no fun! The pilots take turns. Some fly until noon, and some take the afternoon shift. If there are four flights there are four pilots (one taking each flight). And those who did the dropping also took turns. But we (guides) had to stay aboard all day long. That was how it was too with the other Navajos with whom I went to Phoenix.

Some more (Navajos) followed us. They had the same function as we did. They thought it would be easy I heard them saying that. After three days (of flying), a man from Kayenta said that he had been unable to spot his own hogan. He simply couldn't find it. He had been asking that hay be dropped at his own house for three days, but he couldn't find it himself. ``Nope! The plane almost turned over with me (looking for it),'' he said with concern. He told me about this. This man I'm talking about is one called Checked Shirt. When he said the plane nearly turned over with him it was because (this plane was flying low and to gain altitude) the plane had to fly at a steep climb -- (it seemed that it would fall over backward). That's what the pilot (of this plane) told me.

The small planes -- the two motored ones -- loaded at Winslow and hauled from there. These big planes have four motors. They loaded and hauled from Phoenix. They did this because the Winslow Airport was too small for them.

In flight the crew would discuss position. They knew this by means of a map and flight time. They work on a basis of single minutes. They would say it's so far between certain points. And they went on the basis of how many minutes it would take. So this is my experience up above the land of my people. I looked the place over from there.

And it is true. In some places cattle and horses were just standing about, probably on account of the crusted snow. They would break through this crust, and just stand there. They wouldn't go on. Then they would starve or freeze to death right where they stood. In some places the sheep were kept in the corrals all day long. People would merely bring them branches. This is how my people suffered. This was my people's disaster. Many people lost stock. Maybe the range riders and district supervisors feel thankful for it because it reduced the stock.

This is how my people suffered. This is no small matter. God punished us. That's what I think. There is a passage where God says that the people will be punished, and that's what happened. It happened eighteen years ago too. We were punished at that time. We Navajos still haven't found out that God is punishing us. That's the way it is.